\documentclass[11pt]{article}

\usepackage[margin=1in]{geometry}
\usepackage{amsmath,amssymb,amsthm,mathtools}
\usepackage{enumitem}
\usepackage[colorlinks=true,linkcolor=blue,citecolor=blue,urlcolor=blue]{hyperref}

\newtheorem{theorem}{Theorem}
\newtheorem{lemma}[theorem]{Lemma}
\theoremstyle{definition}
\newtheorem{definition}{Definition}[section]
\newtheorem*{remark}{Remark}

\newcommand{\defeq}{\coloneqq}
\newcommand{\Nodes}{N}
\newcommand{\width}{\operatorname{width}}
\newcommand{\children}{\operatorname{children}}
\newcommand{\Occ}{\operatorname{Occ}}

\title{Nice Tree-Decomposition Normalization}
\author{}
\date{}

\begin{document}
\maketitle

\section{Statement}

Graphs are finite simple graphs.  A tree-decomposition of a graph \(G\) is a
pair \((T,\beta)\), where \(T\) is a finite nonempty tree and
\(\beta:\Nodes(T)\to 2^{V(G)}\) assigns a bag to each node, such that:
\begin{enumerate}[label=\textup{(TD\arabic*)},leftmargin=*]
  \item for every \(v\in V(G)\), some node \(t\in\Nodes(T)\) satisfies
    \(v\in\beta(t)\);
  \item for every edge \(uv\in E(G)\), some node \(t\in\Nodes(T)\) satisfies
    \(u,v\in\beta(t)\);
  \item for every \(v\in V(G)\), the set
    \[
      \Occ_T(v)\defeq\{t\in\Nodes(T):v\in\beta(t)\}
    \]
    is connected in \(T\).
\end{enumerate}
The width is at most \(\omega\) if every bag has size at most \(\omega+1\).

\begin{definition}[Nice tree-decomposition]
A rooted tree-decomposition \((T,\beta)\) is \emph{nice} if the following hold.
\begin{enumerate}[label=\textup{(N\arabic*)},leftmargin=*]
  \item The root bag is empty, and every leaf bag is empty.
  \item Every node has at most two children.
  \item If a node \(x\) has exactly one child \(y\), then exactly one vertex is
    changed across the edge \(xy\): either
    \[
      \beta(x)=\beta(y)\cup\{v\}\quad\text{for some }v\notin\beta(y),
    \]
    or
    \[
      \beta(y)=\beta(x)\cup\{v\}\quad\text{for some }v\notin\beta(x).
    \]
    In the first case \(x\) is an introduce node, and in the second case
    \(x\) is a forget node, relative to the top-down orientation.
  \item If a node \(x\) has exactly two children \(y_1,y_2\), then
    \[
      \beta(x)=\beta(y_1)=\beta(y_2).
    \]
    Such a node is a join node.
\end{enumerate}
\end{definition}

\begin{theorem}[Nice normalization]\label{thm:nice-normalization}
Let \(G\) be a finite simple graph.  Let \((T,\beta)\) be a
tree-decomposition of \(G\), and suppose that \(|\beta(t)|\le \omega+1\) for
every \(t\in\Nodes(T)\).  Write \(m=|\Nodes(T)|\).  Then there is a nice
tree-decomposition \((T^\ast,\beta^\ast)\) of \(G\) such that
\[
  |\beta^\ast(x)|\le \omega+1
  \qquad\text{for every }x\in\Nodes(T^\ast),
\]
and
\[
  |\Nodes(T^\ast)|\le (3\omega+5)m.
\]
Consequently, every tree-decomposition of width at most \(\omega\) can be
transformed into a nice tree-decomposition of width at most \(\omega\), with
only \(O_\omega(|\Nodes(T)|)\) nodes.  In particular, if the original width is
exactly \(w\), then taking \(\omega=w\) gives a nice decomposition of width at
most \(w\).
\end{theorem}

\section{Bag-changing paths}

\begin{lemma}[Changing one bag into another]\label{lem:bag-path}
Let \(A,B\subseteq V(G)\) be finite sets with
\(|A|,|B|\le \omega+1\).  There is a rooted path \(P(A,B)\) whose top bag is
\(A\), whose bottom bag is \(B\), and whose consecutive bags differ by exactly
one vertex, except that if \(A=B\) the path consists of a single node.  Every
bag on the path has size at most \(\omega+1\), and the number of edges of the
path is
\[
  |A\setminus B|+|B\setminus A|\le 2(\omega+1).
\]
Moreover, for every vertex \(v\), the nodes of \(P(A,B)\) whose bags contain
\(v\) form a connected subpath, possibly empty.
\end{lemma}

\begin{proof}
Choose enumerations
\[
  A\setminus B=\{a_1,\ldots,a_p\},
  \qquad
  B\setminus A=\{b_1,\ldots,b_q\}.
\]
Start with \(A\).  First delete the vertices \(a_1,\ldots,a_p\), one at a time,
obtaining
\[
  A,\quad
  A\setminus\{a_1\},\quad
  A\setminus\{a_1,a_2\},\quad \ldots,\quad A\cap B.
\]
Then add the vertices \(b_1,\ldots,b_q\), one at a time, obtaining
\[
  A\cap B,\quad
  (A\cap B)\cup\{b_1\},\quad \ldots,\quad B.
\]
This gives a path from \(A\) to \(B\).  Each step changes exactly one vertex,
unless \(p=q=0\), in which case the path has only the single bag \(A=B\).
During the deletion part every bag is a subset of \(A\), and during the
insertion part every bag is a subset of \(B\), so all bags have size at most
\(\omega+1\).  The number of edges is \(p+q\).

For a fixed vertex \(v\), there are four cases.  If \(v\in A\cap B\), then
\(v\) is never deleted and is present throughout the path.  If
\(v\in A\setminus B\), then \(v\) is present until its deletion step and absent
afterwards.  If \(v\in B\setminus A\), then \(v\) is absent until its insertion
step and present afterwards.  If \(v\notin A\cup B\), it is absent throughout.
In every case the occurrence set is a connected subpath or is empty.
\end{proof}

\section{Construction}

Root the original tree \(T\) at an arbitrary node \(r\).  For a node \(t\),
write \(\children(t)\) for its children in this rooted tree, and write
\(T_t\) for the original rooted subtree consisting of \(t\) and all descendants
of \(t\).

We construct, by induction from the leaves upward, a rooted tree \(D_t\) with
bags.  Its root bag will be \(\beta(t)\), every leaf bag will be empty, every
bag will have size at most \(\omega+1\), and all internal nodes will already
satisfy the nice local conditions.

\paragraph{Leaf case.}
If \(t\) has no children, let \(D_t\) be the path
\(P(\beta(t),\emptyset)\).

\paragraph{One-child case.}
Suppose \(t\) has the single child \(s\).  By induction \(D_s\) has root
bag \(\beta(s)\).  Put the path \(P(\beta(t),\beta(s))\) above \(D_s\),
identifying the bottom node of the path with the root of \(D_s\).  If
\(\beta(t)=\beta(s)\), this path has length zero, so no unary node with equal
parent and child bags is inserted.

\paragraph{Many-child case.}
Suppose \(\children(t)=\{s_1,\ldots,s_d\}\) with \(d\ge 2\).  For each
\(i\), put the path \(P(\beta(t),\beta(s_i))\) above \(D_{s_i}\), again
identifying the bottom of the path with the root of \(D_{s_i}\).  The root of
each resulting branch has bag \(\beta(t)\).  Combine these \(d\) branch roots
by a binary join tree whose internal nodes all have bag \(\beta(t)\).  For
instance, use a right-associated tree: the top join node has the first branch
as its left child and a join tree for the remaining branches as its right
child.

After constructing \(D_r\), put the path \(P(\emptyset,\beta(r))\) above it,
identifying the bottom of this path with the root of \(D_r\).  The resulting
rooted tree with bags is denoted \((T^\ast,\beta^\ast)\).

\section{Verification}

\begin{lemma}\label{lem:nice-local}
The pair \((T^\ast,\beta^\ast)\) is nice in the local sense of
Definition~1.1: the root and all leaves have empty bags, every node has at most
two children, unary nodes change one vertex, and binary nodes are join nodes.
All bags have size at most \(\omega+1\).
\end{lemma}

\begin{proof}
The new root is the top of \(P(\emptyset,\beta(r))\), so its bag is empty.
Leaves occur only as the bottom nodes of paths \(P(\beta(t),\emptyset)\) for
original leaves \(t\), and hence have empty bags.

Every path node has at most one child, and by Lemma~\ref{lem:bag-path} adjacent
bags on such a path differ by exactly one vertex.  The only nodes with two
children are the internal nodes of the binary join trees, and by construction
they and their two children all have bag \(\beta(t)\) for the relevant original
node \(t\).  The same lemma gives the width bound for all path nodes; join
nodes have bags equal to original bags, so they also have size at most
\(\omega+1\).
\end{proof}

\begin{lemma}\label{lem:original-bags-survive}
For every original node \(t\in\Nodes(T)\), some node
\(x\in\Nodes(T^\ast)\) satisfies \(\beta^\ast(x)=\beta(t)\).
\end{lemma}

\begin{proof}
We prove the corresponding statement inside each \(D_t\).  If \(t\) is a leaf,
the top node of \(P(\beta(t),\emptyset)\) has bag \(\beta(t)\).  If \(t\) has
one child \(s\) and \(\beta(t)\ne\beta(s)\), the top node of
\(P(\beta(t),\beta(s))\) has bag \(\beta(t)\).  If \(t\) has one child \(s\)
and \(\beta(t)=\beta(s)\), then the induction hypothesis for \(s\) gives a
node with bag \(\beta(s)=\beta(t)\).  If \(t\) has at least two children, the
root of the binary join tree has bag \(\beta(t)\).
\end{proof}

\begin{lemma}\label{lem:decomp-axioms}
The pair \((T^\ast,\beta^\ast)\) is a tree-decomposition of \(G\).
\end{lemma}

\begin{proof}
The vertex-covering axiom follows from Lemma~\ref{lem:original-bags-survive}
and the vertex-covering axiom for \((T,\beta)\).  Indeed, if \(v\in V(G)\),
choose \(t\in\Nodes(T)\) with \(v\in\beta(t)\); then a node of \(T^\ast\) has
bag \(\beta(t)\), hence contains \(v\).  The edge-covering axiom is identical:
if \(uv\in E(G)\), choose \(t\) with \(u,v\in\beta(t)\), and use the surviving
copy of \(\beta(t)\).

It remains to prove the running-intersection axiom.  Fix \(v\in V(G)\).  Let
\[
  S_v\defeq\{t\in\Nodes(T):v\in\beta(t)\}.
\]
This set is connected in the original tree \(T\).  We compare the occurrence
set of \(v\) in \(T^\ast\) with \(S_v\).

Each added path \(P(\beta(t),\beta(s))\) above a child \(s\) contributes a
connected subpath of nodes containing \(v\), by Lemma~\ref{lem:bag-path}.  If
\(v\) belongs to both endpoint bags, this subpath connects the side attached to
\(t\) with the side attached to \(s\).  If \(v\) belongs to exactly one endpoint
bag, the subpath is attached only at that endpoint.  If \(v\) belongs to
neither endpoint bag, the path contributes nothing.  The same description holds
for the top path \(P(\emptyset,\beta(r))\) and for the leaf paths
\(P(\beta(t),\emptyset)\).

Every join tree introduced at an original node \(t\) has all bags equal to
\(\beta(t)\).  Hence it is entirely contained in the occurrence set of \(v\)
when \(t\in S_v\), and it is disjoint from that occurrence set when
\(t\notin S_v\).

Thus the occurrence set of \(v\) in \(T^\ast\) is obtained from the connected
original subtree \(S_v\) by replacing original nodes and original edges with
connected pieces, and by attaching additional connected initial or terminal
path segments at nodes of \(S_v\).  Such replacements and attachments preserve
connectedness.  Therefore
\[
  \{x\in\Nodes(T^\ast):v\in\beta^\ast(x)\}
\]
is connected.
\end{proof}

\begin{proof}[Proof of Theorem~\ref{thm:nice-normalization}]
By Lemmas~\ref{lem:nice-local} and~\ref{lem:decomp-axioms},
\((T^\ast,\beta^\ast)\) is a nice tree-decomposition of \(G\) with all bags of
size at most \(\omega+1\).  It remains only to record the size bound.

Let \(m=|\Nodes(T)|\).  Each original edge \(ts\), oriented from parent \(t\)
to child \(s\), contributes at most \(2(\omega+1)\) new path nodes above
\(D_s\), because the bottom node is identified with the root of \(D_s\).
There are \(m-1\) original edges.  Each original leaf \(t\) contributes at most
\(\omega+2\) nodes in the terminal path \(P(\beta(t),\emptyset)\).  There are
at most \(m\) original leaves.  The binary join trees contribute
\(\max(d(t)-1,0)\) join nodes at a node with \(d(t)\) children, and
\[
  \sum_{t\in\Nodes(T)}\max(d(t)-1,0)
  \le \sum_{t\in\Nodes(T)} d(t)
  =m-1.
\]
Finally, the top path \(P(\emptyset,\beta(r))\) contributes at most
\(\omega+1\) new nodes above \(D_r\).

Therefore
\[
\begin{aligned}
  |\Nodes(T^\ast)|
  &\le 2(\omega+1)(m-1)+(\omega+2)m+(m-1)+(\omega+1) \\
  &= (3\omega+5)m-\omega-2 \\
  &\le (3\omega+5)m.
\end{aligned}
\]
This proves the claimed bound.
\end{proof}

\begin{remark}
The constant \(3\omega+5\) is not intended to be optimal.  The important point
for Courcelle's theorem is that the width is not increased and the number of
nodes is linear in the size of the supplied decomposition for fixed
\(\omega\).
\end{remark}

\end{document}
